% Scope (SPEC table of contents): Abbreviations the book uses. Grows with the chapters.
%
% An entry here needs an expansion verified against the source, not remembered.
% ERASER is deliberately absent for that reason: its own title does not expand
% the acronym and chapter 08's session found no source that does, so it sits in
% appendix B's keep-in-English block as a benchmark name instead of being given
% an invented expansion here.
%
% SCM is absent on the same rule, found by chapter 09's session: paper 21 uses
% the abbreviation across all 31 pages and expands it nowhere, so it too sits in
% appendix B as a bare metric name. FUR is present because that paper does
% expand it (p. 24). AUROC is present by a different route, since paper 21 uses
% it unexpanded as well: its confidence intervals are DeLong intervals, and the
% DeLong paper it cites for them supplies "receiver operating characteristic".
% research/ch09-bai-bao-neo.md records that chain.
%
% NIS is absent on the CQS rule, found by chapter 15's session: its source does
% expand it, but the chapter names it once and then never uses the abbreviation
% again, so it sits in appendix B's keep-in-English block instead. That is the
% third time that rule has cut, after CQS and ASV.
%
% CBM was seeded at init and sat unused for fourteen chapters. Chapter 15
% earns it, on the precedent chapter 13 set with AIT: a session that touches
% the subject of a seeded abbreviation decides whether to use it or remove the
% row, and this one introduces CBM at first use in section 15.1 and then uses
% it throughout.
\chapter{Bảng chữ viết tắt}
\label{app:viet-tat}

\begin{tabular}{@{}lp{9cm}@{}}
\toprule
Viết tắt & Nghĩa \\
\midrule
XAI & explainable AI, trí tuệ nhân tạo giải thích được \\
LIME & local interpretable model-agnostic explanations \\
SHAP & SHapley Additive exPlanations \\
IG & integrated gradients \\
Grad-CAM & gradient-weighted class activation mapping \\
CBM & concept bottleneck model \\
CoT & chain of thought, chuỗi suy luận \\
LLM & large language model, mô hình ngôn ngữ lớn \\
RLHF & reinforcement learning from human feedback \\
SAE & sparse autoencoder \\
AIT & algorithmic information theory \\
RISE & randomized input sampling for explanation \\
ROAR & remove and retrain \\
AUROC & area under the receiver operating characteristic curve \\
FUR & faithfulness via unlearning reasoning steps \\
STII & Shapley-Taylor interaction index \\
SOP & sum of power \\
FGSM & fast gradient sign method \\
PC & Peter-Clark, tên thuật toán khám phá nhân quả \\
IDA & intervention calculus when the DAG is absent \\
RMSE & root mean squared error \\
CEM & concept embedding model \\
CTL & concepts-task leakage \\
ICL & interconcept leakage \\
OIS & oracle impurity score \\
CB-SAE & concept bottleneck sparse autoencoder \\
\bottomrule
\end{tabular}
